\subsection{Discrete measures}

PROPOSITION 14. --- For a measure \(\mu\) on a locally compact space \(X\) to be a discrete measure (§1, No.~3, Example I), it is necessary and sufficient that \(\operatorname{Supp}(\mu)\) be a discrete closed subspace of \(X\) .

Let \(\mu\) be a discrete measure on X, defined by the masses \(h(x) \neq 0\) placed at the points \(x\) of a discrete closed subspace N of X; let us show that \(\operatorname{Supp}(\mu) = \mathbf{N}\) . For every \(a \in \mathbf{N}\) and every neighborhood V of \(a\) , there exists a function \(f \in \mathcal{K}(\mathbf{X}; \mathbf{C})\) with support contained in V, equal to 1 at the point \(a\) and to 0 at the other points of N, whence \(\mu(f) = h(a) \neq 0\) . On the other hand if \(b \notin \mathbf{N}\) then there exists a neighborhood W of \(b\) not intersecting N; for every function \(g \in \mathcal{K}(\mathbf{X}; \mathbf{C})\) with support contained in W, we therefore have \(\mu(g) = 0\) , which proves that \(b \notin \operatorname{Supp}(\mu)\) .

Conversely, let \(\mu\) be a measure such that \(\operatorname{Supp}(\mu)\) is a discrete closed subspace N of X. For every \(a \in N\) , there exists an open neighborhood \(V_{a}\) of a that contains no point of N distinct from a; the restriction of \(\mu\) to \(V_{a}\) is therefore a point measure with support \(\{a\}\) (No.~2, Prop. 5), hence (No.~4, Prop. 12) is of the form \(h(a)\varepsilon_{a}\) , where \(h(a)\neq0\) . Setting \(h(x)=0\) at the points of \(\complement N\), and denoting by \(\nu\) the measure defined by the masses \(h(x)\) , the principle of localization shows that \(\nu=\mu\) .

We thus see that on a discrete space X, every measure is discrete.

